\subsubsection{Variability in Regression Reduction}
\label{subsubsec:regression-reduction-violin}

Figure~\ref{fig:regression-reduction-violin} presents violin plots that reveal the complete probability density of regression reduction outcomes. The width at any vertical position represents probability density; wider sections indicate common outcomes, narrower sections indicate rare outcomes. Red and blue lines mark mean and median respectively.

The violin shapes reveal a stark categorical divide across coding agents. Claude Code achieves 93.0\% positive outcomes (346 of 372 repairs), 5.6\% zero-change (21 repairs), and only 1.3\% negative (5 repairs). OpenAI Codex shows similar behavior with 89.9\% positive, 7.9\% zero, and 2.2\% negative. Both exhibit narrow distributions with low variance ($\sigma=5.91$ and $6.64$), indicating predictable behavior. Their worst-case scenarios are bounded (min $-18$ and $-47$). In contrast, Qwen Code achieves only 30.2\% positive outcomes (110 of 364 repairs), with 59.9\% zero-change (218 repairs compile but produce no improvement) and 9.9\% negative (36 repairs introduce regressions). Gemini CLI shows 46.1\% positive, 43.5\% zero, and 10.4\% negative. Both display wide, dispersed distributions with high variance ($\sigma=28.57$ and $31.25$), nearly 5$\times$ larger than Claude/Codex—and catastrophic worst-cases (min $-543$ and $-488$) where a single repair break hundreds of tests.


The variance ratio ($\sigma_{\text{Gemini}}^2 / \sigma_{\text{Claude}}^2 \approx 28$) quantifies operational risk: high-variance tools require substantially more quality assurance resources per repair. Distribution shape predicts automation feasibility—narrow violins enable autonomous deployment in CI/CD pipelines, while wide violins necessitate supervised deployment with mandatory human review. Gemini's range $[-488, +26]$ includes catastrophic outcomes that could halt development, whereas Claude's $[-18, +73]$ contains only manageable scenarios. This tail risk asymmetry explains why practitioners need to be careful while leveraging high-variance coding assistance tools despite seemingly acceptable averages. When selecting tools to build production-ready applications, practitioners must consider different aspects: central tendency (mean, median) indicates typical outcomes, dispersion (standard deviation) reveals consistency, and worst-case tail risk (minimum/maximum values) exposes catastrophic failures.

\begin{figure}[t]
\centering
\includegraphics[width=\columnwidth]{analysis/figures/agent_regression_reduction_violin.pdf}
\caption{Violin plot distributions of regression reduction across four agentic coding tools (372 Defects4J bugs). Width represents probability density; red/blue lines mark mean/median; black lines span observed range. Claude Code ($\mu=2.47$, $\sigma=5.91$) and OpenAI Codex ($\mu=2.25$, $\sigma=6.64$) exhibit narrow distributions with 93.0\% and 89.9\% positive outcomes, indicating production-ready consistency. Qwen Code ($\mu=-1.34$, $\sigma=28.57$) and Gemini CLI ($\mu=-1.92$, $\sigma=31.25$) exhibit wide distributions with 59.9\% and 43.5\% zero-change outcomes and catastrophic minimums of $-543$ and $-488$, indicating high-variance behavior requiring extensive human oversight. The visualization demonstrates that distribution shape—not mean alone—determines practical utility.}
\label{fig:regression-reduction-violin}
\end{figure}
